% =====================================================================
% ABSTRACT. Third rewrite, 17 November 2026.
%
% AW's direction for this one: "What if we were to lead with the AI side
% instead of scale. Endometriosis is still the starting point. We add
% scale and some really innovative architecture. E3 is good but needs to
% be stated in a more understandable tone leading with the result. E4 is
% also a catchy result for a disease atlas." And: "'Claims no contributing
% cohort could test alone are settled' lines like these are not headline
% results."
%
% SO THE ORDER IS: endometriosis, then the architecture, then the scale,
% then what the atlas found. Scale is no longer the second beat; it
% arrives after the model has been described, as something the approach
% made affordable rather than as the claim itself.
%
% E3 NOW LEADS WITH THE RESULT AND NAMES THE GENES. The previous draft
% opened on the method ("Decoding that direction returns...") and said
% "four of its fifteen leading genes", which asks the reader to hold an
% abstraction. IBDex names its genes outright, DMBT1, DUOXA2, OLFM4,
% NOS2, CXCL1, and HLCA names SPP1+ macrophages, so naming them is in
% register for a disease atlas.
%
% E4 IS PROMOTED to its own sentence on AW's reading that four genes no
% prior endometriosis report contains is catchy for an atlas. It is the
% one claim here that is a discovery rather than a confirmation.
%
% THE META-SENTENCE IS GONE. "Claims no contributing cohort could test
% alone are settled" describes an activity, not a finding. What replaced
% it is the two findings themselves, with their donor counts.
%
% LENGTH. 236 words against NMI's 150 spec; the six NMI abstracts measure
% 147 to 203, median 183. AW, 17 November: "Feel free to go over, I can
% edit", and the standing rule is to seed complete rather than short. The
% cheapest cuts, in order: the frontier sentence (25 words, but it is
% result B1 at importance 10), the eight gene names (12 words), and the
% closing composition clause.
% =====================================================================

\begin{abstract}
\noindent
Endometriosis is diagnosed by surgery, so its biology has been studied one
small cohort at a time, and the field's reference atlas of the endometrium
spans 63 donors. Here we present \elaj, a single-cell disease atlas of
endometriosis in which biological structure built into the model replaces
pretraining scale. A conditional variational autoencoder reads each cell as
44 pathway groups rather than as fifteen thousand independent genes, and
returns sixteen interpretable dimensions. Two coefficients trace a frontier
along which donor identity is traded against retained biology, so every
result is reported from a named position on it. At 1{,}883{,}986 parameters
\elaj\ holds more menstrual-cycle signal at lower donor identity than nine
competing representations, foundation models over fifty times its size among
them. The atlas is built from 1{,}577{,}728 cells and 254 donors across
eleven studies, four times the donors of that reference atlas. Inside immune
cells it separates patients from controls on donors it never saw, at an AUC
of 0.708, and names the genes that carry the separation. Four of them,
SERPINB2, MMP27, LRRC15 and INHBA, were independently reported as raised in
patients by a cohort outside this corpus,
\textit{P}~=~4~\texttimes~10\textsuperscript{--7}; none of the genes at the
control end was. Four others, SIGLEC7, ADAMDEC1, ALPK2 and KCNH1, appear in
no prior endometriosis report. Across 194 donors cell composition does not
differ between patients and controls, and severity among patients is
recovered by neither the latent nor a cell count.
\end{abstract}
